\subsection{Ethical Considerations in AI Research and Experimentation}
\label{sec:2.4.2}
AI research and experimentation raise ethical stakes in proportion to how far a system's cognition might extend toward the awareness section~\ref{sec:2.4.1} defines. Six principles organize what follows, and three of them do work here that they do nowhere else in this book.

\runin{Beneficence}
Beneficence commits researchers to promoting well-being and avoiding harm — to the humans an AI system serves and, where the system is a candidate sentient system in the sense of section~\ref{sec:2.4.1}, to the system itself. Decisions about programming, reprogramming, or terminating such a system carry ethical weight in proportion to where it sits on the nociception-to-suffering ladder, and not to how advanced its behavior looks from outside.

\runin{The Precautionary Principle}
Precaution asks for risk assessment and safeguards before deployment rather than after. Applied to candidate sentient systems it also covers the psychological and societal fallout of introducing artificial entities capable of experience — effects on human behavior and social norms that regulatory frameworks built for non-sentient AI assumed away rather than addressed, and that call for guidelines drafted with the difference in mind.

\runin{Interdisciplinary Review and Oversight}
The ethical questions AI research raises span expertise no single field holds, and concentrating that judgment in any one of them repeats, in miniature, the concentration-of-power problem this book treats as a harm on its own terms. Independent ethical review boards, staffed across disciplines, are the mechanism for turning the principle into practice: reviewing a proposed study before it begins, inspecting it as it runs, and evaluating it against measurable criteria once results are in, on the model medical institutional review boards already provide for drug trials.

The remaining three bear on this research the way they bear on everything else the field produces. Autonomy asks whether a candidate sentient system can affect decisions about its own operation and welfare, which section~\ref{sec:2.4.4} takes up as the prior question of whether consent is available to such a system at all. Justice asks whether an AI system's benefits and burdens fall evenly across the people it touches, which is chapter~\ref{sec:6}'s subject. And transparency asks that a study be documented well enough for peers to reproduce it and that a deployed system be able to account for its own outputs precisely enough for a person to contest them, which is chapter~\ref{sec:9}'s.
